\documentclass[11pt,a4paper]{article}
\usepackage[margin=1in]{geometry}
\usepackage[T1]{fontenc}
\usepackage[utf8]{inputenc}
\usepackage{amsmath,amssymb,booktabs,longtable,array}
\usepackage{xcolor}
\usepackage[hidelinks]{hyperref}
\newcolumntype{L}[1]{>{\raggedright\arraybackslash}p{#1}}
\newcommand{\code}[1]{\begingroup\urlstyle{tt}\footnotesize\nolinkurl{#1}\endgroup}
\setlength{\LTleft}{0pt}
\setlength{\LTright}{0pt}
\setlength{\tabcolsep}{4pt}
\setlength{\emergencystretch}{3em}
\allowdisplaybreaks
\renewcommand{\arraystretch}{1.2}

\title{Particle-Based Export Formula Report}
\author{Generated from stress\_particle\_based.py}
\date{}
\begin{document}
\maketitle
\noindent Formulas implemented in \code{stress_particle_based.py}\par
\medskip
\noindent This document summarizes the exact formulas currently implemented in \code{stress_particle_based.py} to produce exported CSV data.

\section*{Output files}
\begin{longtable}{@{}L{0.18\textwidth}L{0.32\textwidth}L{0.42\textwidth}@{}}
\toprule
\textbf{Type} & \textbf{File name pattern} & \textbf{Content}\\
\midrule
\endhead
Summary CSV & \code{<dem>_particle_over_time.csv} & One row per processed timestep.\\
Extracted CSV & \code{<dem>_particle_extracted_at_<time>.csv} & Per-particle values at the selected extract time.\\
\bottomrule
\end{longtable}

\section*{Notation}
\begin{longtable}{@{}L{0.22\textwidth}L{0.70\textwidth}@{}}
\toprule
\textbf{Symbol} & \textbf{Meaning}\\
\midrule
\endhead
$k$ & Current timestep index.\\
$t_k$ & Current sample time returned by EDEM for timestep k.\\
$\Delta t_k$ & Interval length from the previous sample to the current sample.\\
$C_k$ & Set of collisions present in the current surfSurf sample.\\
$N_k$ & Number of sampled collisions at timestep k.\\
$N_k^{new}$ & Number of sampled collisions whose start time lies inside the current interval.\\
$p \in P$ & Particle index in the cached particle metadata used by the script.\\
$i \sim p$ & Collision i touches particle p as either first or second particle.\\
$I[\cdot]$ & Indicator function equal to 1 when the stated condition is true and 0 otherwise.\\
$n_{p,k}$ & Number of sampled current collisions touching particle p at timestep k.\\
$m_p, r_p$ & Mass and radius of particle p from the cached particle metadata.\\
$s_i, e_i$ & Start time and end time of current collision i.\\
$d_i$ & Full contact duration of collision i over its whole lifetime.\\
$E_{n,i}, E_{s,i}, E_i$ & Collision energies used in the code.\\
$F_{n,i}, F_{s,i}, F_i$ & Current force magnitudes used in the code.\\
$F_i^{max}$ & Maximum total collision-force magnitude used for impact filtering.\\
$R_k$ & Previous-only collisions stored in state that disappeared before the current sample but still overlap the interval.\\
$\tau_i^{overlap}$ & Overlap duration of a current collision with the interval $[t_{k-1}, t_k]$.\\
$\tau_j^{tail}$ & Overlap duration attributed to a previous-only tail collision $j$ in $R_k$.\\
$J_{n,p,k}, J_{s,p,k}, J_{p,k}$ & Normal, shear, and total particle impulse attributed to the interval.\\
$T_{p,k}^{contact}$ & Total contact time assigned to particle p in the interval.\\
$A_p$ & Particle area proxy used in the stress calculation, computed from radius.\\
$\sigma_{n,p}, \sigma_{s,p}, \sigma_p$ & Exported normal, shear, and total stress-like values per particle.\\
$\boldsymbol{\omega}_h$ & Angular-velocity vector of host particle h.\\
$d_{h,k}^{rot}, D_{h,k}^{rot}$ & Current and cumulative rotational travel distance for host particle h.\\
$\xi_{h,k}$ & Binary indicator showing whether host particle h is contacting at timestep k.\\
$T_{h,k}^{contact}$ & Cumulative time for which host particle h is marked as contacting.\\
$H_k$ & Host particles selected by the configured host material.\\
\bottomrule
\end{longtable}

\section*{How to read the symbols}
\paragraph{Time and interval symbols}
$k$ is only a sample index, while $t_k$ is the physical simulation time at that sample.
$\Delta t_k$ is the width of the reporting interval, so any rate, frequency, or power term is normalized by this interval length.
For a collision that lives longer than one reporting interval, $\tau_i^{overlap}$ keeps only the part that lies inside $[t_{k-1}, t_k]$, and $\tau_j^{tail}$ keeps the leftover tail contribution from a collision that disappeared before the current sample but still overlaps the interval boundary.
So $\tau_i^{overlap}$ is an interval-only portion of a current collision, whereas $\tau_j^{tail}$ is a carry-over portion from a previous-only collision state.

\paragraph{Set and counting symbols}
$C_k$ is the full set of collisions that are visible in the current surfSurf sample, not only the newly started ones.
$R_k$ is not a new collision set from EDEM and not part of the current sampled set $C_k$; it is an internal carry-over state used so the particle-level impulse does not lose the tail end of a contact between samples.
$P$ is the cached particle population, $H_k$ is the host-material subset of interest, and $i \sim p$ simply means collision $i$ touches particle $p$.
The indicator $I[\cdot]$ turns a logical statement into a numeric 0 or 1, which is why it appears in counting formulas such as $n_{p,k}$.

\paragraph{Mechanical symbols}
$m_p$ and $r_p$ are particle properties, $A_p = \pi r_p^2$ is the area proxy built from radius, $E$ terms are energies, $F$ terms are force magnitudes, and $J$ terms are impulses obtained by integrating force over time.
A bar, as in $\bar{F}_{p,k}$, denotes an average over the current set of particle-touching collisions.
$T_{p,k}^{contact}$ is the total contact time assigned to a particle inside the interval, and dividing impulse by contact time and area produces the exported stress-like values $\sigma_{n,p}$, $\sigma_{s,p}$, and $\sigma_p$.
These $\sigma$ values are code-defined stress proxies, not direct continuum stresses or true local contact stresses.

\paragraph{Host-particle rotation symbols}
For a host particle $h$, $\boldsymbol{\omega}_h$ is angular velocity, $d_{h,k}^{rot}$ is the rotational distance traveled during the current interval, and $\xi_{h,k}$ is a contact flag that decides whether this distance contributes to the cumulative value $D_{h,k}^{rot}$.
$T_{h,k}^{contact}$ accumulates the amount of interval time during which the host particle is considered to be in contact.

\section*{Core equations}
\paragraph{P1. Interval length and overlap duration}
\[
\begin{aligned}
\Delta t_k &= \max(0, t_k - t_{k-1}), \\
\tau_i^{overlap} &= \max\!\left(0, \min(e_i, t_k) - \max(s_i, t_{k-1})\right)
\end{aligned}
\]
The particle report inherits the same interval model as the collision report.

\paragraph{P2. Current collision count per particle}
\[
n_{p,k} = \sum_{i \in C_k} I[i \sim p]
\]
This count uses the current sampled collisions at timestep k, not unique collisions over the entire interval.

\paragraph{P3. Average current collision-force magnitude per particle}
\[
\bar{F}_{p,k} = \frac{\sum_{i \in C_k,\ i \sim p} F_i}{n_{p,k}}
\]
If $n_{p,k}$ is zero, the code writes zero for that particle.

\paragraph{P4. Energy split from collisions to particles}
\[
\begin{aligned}
E_{n,p} &= \frac{1}{2}\sum_{i \in C_k,\ i \sim p} E_{n,i}, \\
E_{s,p} &= \frac{1}{2}\sum_{i \in C_k,\ i \sim p} E_{s,i}, \\
E_p &= \frac{1}{2}\sum_{i \in C_k,\ i \sim p} E_i
\end{aligned}
\]
Each collision contributes half of its energy to each particle in the pair.

\paragraph{P5. Particle specific-intensity values}
\[
SI_{n,p} = \frac{E_{n,p}}{m_p},\quad SI_{s,p} = \frac{E_{s,p}}{m_p},\quad SI_p = \frac{E_p}{m_p}
\]
If particle mass is zero, the code writes zero.

\paragraph{P6. Aggregate specific energy by particle}
\[
\begin{aligned}
SE_{n,k}^{part} &= \frac{\sum_{p \in P} E_{n,p}}{\sum_{p \in P} m_p}, \\
SE_{s,k}^{part} &= \frac{\sum_{p \in P} E_{s,p}}{\sum_{p \in P} m_p}, \\
SE_k^{part} &= \frac{\sum_{p \in P} E_p}{\sum_{p \in P} m_p}
\end{aligned}
\]
The implementation aggregates particle energies first, then divides by the sum of valid particle masses.
So this aggregate specific energy is not the same as taking the arithmetic mean of the per-particle $SI$ values.

\paragraph{P7. Trapezoidal impulse from current matched collisions}
\[
J_{n,p,k}^{curr} = \sum_{i \in C_k,\ i \sim p} \frac{1}{2}\left(F_{n,i}^{prev} + F_{n,i}^{curr}\right)\tau_i^{overlap}
\]
The same trapezoidal matching rule as the collision report is used for current matched collisions.
Only collisions that are present in the current sampled set $C_k$ and matched to previous state contribute to this current term.

\paragraph{P8. Previous-only tail impulse added only at particle level}
\[
\widetilde{J}_{n,p,k}^{tail} = \sum_{j \in R_k,\ j \sim p} \frac{1}{2}F_{n,j}^{prev}\tau_j^{tail}
\]
Equivalent tail terms are also used for tangential and total impulse. This term is not added to collision-level exports.
$R_k$ here is the previous-only carry-over state, so these tail terms can contribute even when the corresponding contact is no longer present in the current sampled set $C_k$.

\paragraph{P9. Total particle impulse and accumulated contact time}
\[
\begin{aligned}
J_{n,p,k} &= J_{n,p,k}^{curr} + \widetilde{J}_{n,p,k}^{tail}, \\
T_{p,k}^{contact} &= \sum_{i \in C_k,\ i \sim p}\tau_i^{overlap} + \sum_{j \in R_k,\ j \sim p}\tau_j^{tail}
\end{aligned}
\]
The same structure is used for shear and total impulse.
Therefore particle impulse and particle contact time may include both currently sampled contacts and previous-only tail contributions.

\paragraph{P10. Particle area and stress values}
\[
\begin{aligned}
A_p &= \pi r_p^2, \\
\sigma_{n,p} &= \frac{J_{n,p,k}}{T_{p,k}^{contact} A_p}, \\
\sigma_{s,p} &= \frac{J_{s,p,k}}{T_{p,k}^{contact} A_p}, \\
\sigma_p &= \frac{J_{p,k}}{T_{p,k}^{contact} A_p}
\end{aligned}
\]
If contact time or particle area is zero, the code writes zero.
These exported $\sigma$ values should therefore be read as stress-like proxies from the implementation.

\paragraph{P11. Particle power values}
\[
\begin{aligned}
P_{n,p} &= \frac{E_{n,p}}{\Delta t_k}, \\
P_{s,p} &= \frac{E_{s,p}}{\Delta t_k}, \\
P_p &= \frac{E_p}{\Delta t_k}
\end{aligned}
\]
Power is computed from the current interval length.

\paragraph{P12. Host-particle rotational metrics}
\[
\begin{aligned}
d_{h,k}^{rot} &= r_h \|\boldsymbol{\omega}_h\| \Delta t_k, \\
\xi_{h,k} &= I[n_{h,k} > 0], \\
D_{h,k}^{rot} &= D_{h,k-1}^{rot} + \xi_{h,k} d_{h,k}^{rot}, \\
T_{h,k}^{contact} &= T_{h,k-1}^{contact} + \xi_{h,k}\Delta t_k
\end{aligned}
\]
The contact indicator is based on the current sampled active collision count for each host particle.
It is a gate based only on whether the host particle has any current sampled collision at timestep $k$; it is not itself an accumulated overlap measure.

\paragraph{P13. Collision-side fields still exported in particle\_over\_time.csv}
\[
\begin{aligned}
\overline{d}_k^{coll} &= \frac{1}{N_k}\sum_{i \in C_k} d_i, \\
\overline{\tau}_k^{coll} &= \frac{1}{N_k}\sum_{i \in C_k}\tau_i^{overlap}, \\
f_k^{coll} &= \frac{N_k^{new}}{\Delta t_k}
\end{aligned}
\]
The particle summary CSV also keeps several collision-derived summary columns from the shared timestep core.
In particular, the $f_k^{coll}$ term here uses the same new-collision subset rule as the collision report, while $C_k$ itself remains the full current sampled collision set.

\section*{Summary CSV fields}
\small
\begin{longtable}{@{}L{0.33\textwidth}L{0.22\textwidth}L{0.37\textwidth}@{}}
\toprule
\textbf{Field} & \textbf{Calculation / reference} & \textbf{Meaning}\\
\midrule
\endhead
\code{time} & Current sample time $t_k$. & EDEM time value for the processed timestep.\\
\code{number_of_collisions} & $N_k = |C_k|$. & Number of sampled collisions in the current surfSurf snapshot.\\
\code{Ave_number_of_collisions_per_particle} & Mean over particles of $n_{p,k}$. Eq.~P2. & Average number of sampled current collisions per particle.\\
\code{Ave_contact_duration_per_collision} & Mean of $d_i$ over sampled collisions. Eq.~P13. & Average full contact duration of current sampled collisions.\\
\code{total_collision_time} & Sum of $\tau_i^{overlap}$ over sampled collisions. Eq.~P13. & Total overlap time attributed to the interval.\\
\code{average_collision_time} & Mean of $\tau_i^{overlap}$ over sampled collisions. Eq.~P13. & Average overlap time per sampled collision in the current interval.\\
\code{collision_frequency} & $N_k^{new} / \Delta t_k$. Eq.~P13. & Frequency of newly started sampled collisions in the interval.\\
\code{Ave_TotalColForce_inSys_byParticle} & Mean over particles of $\bar{F}_{p,k}$. Eq.~P3. & Average current collision-force magnitude per particle, including zero for particles with no current collision.\\
\code{high_impact_event_fraction} & Count($F_i^{max} > \mathrm{threshold}$) / $N_k$ using collision max-force magnitude. & Fraction of sampled collisions above the configured impact threshold.\\
\code{Ave_SI_normal_per_particle} & Mean over particles of $SI_{n,p}$. Eq.~P5. & Average normal specific intensity per particle.\\
\code{Ave_SI_shear_per_particle} & Mean over particles of $SI_{s,p}$. Eq.~P5. & Average shear specific intensity per particle.\\
\code{Ave_SI_total_per_particle} & Mean over particles of $SI_p$. Eq.~P5. & Average total specific intensity per particle.\\
\code{Total_normal_specific_energy_by_particle} & Aggregate particle normal specific energy. Eq. P6. & Normal energy summed over particles, then divided by summed particle mass.\\
\code{Total_shear_specific_energy_by_particle} & Aggregate particle shear specific energy. Eq. P6. & Shear energy summed over particles, then divided by summed particle mass.\\
\code{Total_specific_energy_by_particle} & Aggregate particle total specific energy. Eq. P6. & Total energy summed over particles, then divided by summed particle mass.\\
\code{Ave_normal_stress_value_per_particle} & Mean over particles of $\sigma_{n,p}$. Eq.~P10. & Average normal stress proxy per particle.\\
\code{Ave_shear_stress_value_per_particle} & Mean over particles of $\sigma_{s,p}$. Eq.~P10. & Average shear stress proxy per particle.\\
\code{Ave_stress_value_per_particle} & Mean over particles of $\sigma_p$. Eq.~P10. & Average total stress proxy per particle.\\
\code{Ave_normal_power_per_particle} & Mean over particles of $P_{n,p}$. Eq.~P11. & Average normal power per particle.\\
\code{Ave_shear_power_per_particle} & Mean over particles of $P_{s,p}$. Eq.~P11. & Average shear power per particle.\\
\code{Ave_power_per_particle} & Mean over particles of $P_p$. Eq.~P11. & Average total power per particle.\\
\code{Total_normal_power_by_particle} & Sum over particles of $P_{n,p}$. Eq.~P11. & Total normal power across particles.\\
\code{Total_shear_power_by_particle} & Sum over particles of $P_{s,p}$. Eq.~P11. & Total shear power across particles.\\
\code{Total_power_by_particle} & Sum over particles of $P_p$. Eq.~P11. & Total power across particles.\\
\code{Ave_rotational_distance_per_host_particle} & Mean of $d_{h,k}^{rot}$ over all host particles sampled at $k$. Eq.~P12. & Average host rotational distance in the current interval.\\
\code{Ave_rotational_distance_total_per_host_particle} & Mean of $D_{h,k}^{rot}$ over host particles with $\xi_{h,k} = 1$. Eq.~P12. & Average cumulative host rotational distance among currently contacting host particles.\\
\code{Ave_cumulative_contact_time_per_host_particle} & Mean of $T_{h,k}^{contact}$ over host particles with $\xi_{h,k} = 1$. Eq.~P12. & Average cumulative host contact time among currently contacting host particles.\\
\bottomrule
\end{longtable}
\normalsize

\section*{Extracted CSV fields}
\small
\begin{longtable}{@{}L{0.33\textwidth}L{0.22\textwidth}L{0.37\textwidth}@{}}
\toprule
\textbf{Field} & \textbf{Calculation / reference} & \textbf{Meaning}\\
\midrule
\endhead
\code{time} & Scalar $t_k$ for the extract timestep. & Written once at the top of the extracted block by the CSV expansion helper.\\
\code{particle_id} & Cached particle ID from the particle metadata. & Particle identifier used by the script.\\
\code{particle_mass} & Cached mass $m_p$ from the particle metadata. & Particle mass used in particle specific-intensity calculations.\\
\code{collision_number_per_particle} & $n_{p,k}$ from Eq.~P2. & Current sampled collision count touching the particle.\\
\code{Ave_TotalColForce_per_particle} & $\bar{F}_{p,k}$ from Eq.~P3. & Average current collision-force magnitude touching the particle.\\
\code{SI_normal_per_particle} & $SI_{n,p}$ from Eq.~P5. & Normal specific intensity per particle.\\
\code{SI_shear_per_particle} & $SI_{s,p}$ from Eq.~P5. & Shear specific intensity per particle.\\
\code{SI_total_per_particle} & $SI_p$ from Eq.~P5. & Total specific intensity per particle.\\
\code{specific_energy_normal_per_particle} & Alias of \code{SI_normal_per_particle} in the code. & Same numerical value as \code{SI_normal_per_particle}.\\
\code{specific_energy_shear_per_particle} & Alias of \code{SI_shear_per_particle} in the code. & Same numerical value as \code{SI_shear_per_particle}.\\
\code{specific_energy_total_per_particle} & Alias of \code{SI_total_per_particle} in the code. & Same numerical value as \code{SI_total_per_particle}.\\
\code{normal_stress_value_per_particle} & $\sigma_{n,p}$ from Eq.~P10. & Normal stress proxy.\\
\code{shear_stress_value_per_particle} & $\sigma_{s,p}$ from Eq.~P10. & Shear stress proxy.\\
\code{stress_value_per_particle} & $\sigma_p$ from Eq.~P10. & Total stress proxy.\\
\code{normal_power_per_particle} & $P_{n,p}$ from Eq.~P11. & Normal power.\\
\code{shear_power_per_particle} & $P_{s,p}$ from Eq.~P11. & Shear power.\\
\code{power_per_particle} & $P_p$ from Eq.~P11. & Total power.\\
\code{rotational_distance_contact_xi_per_particle} & $\xi_{h,k}$ from Eq.~P12 for host particles, else $0$. & Current contact indicator written into the particle table.\\
\code{rotational_distance_per_particle} & $d_{h,k}^{rot}$ from Eq.~P12 for host particles, else $0$. & Current rotational distance for host particles.\\
\code{rotational_distance_total_per_particle} & $D_{h,k}^{rot}$ from Eq.~P12 for host particles, else $0$. & Cumulative rotational distance for host particles while in contact.\\
\bottomrule
\end{longtable}
\normalsize

\section*{Implementation notes}
\noindent - All displayed formulas reflect the current implementation in \code{stress_particle_based.py}.\par
\noindent - Particle-level impulse uses both the matched current trapezoid term and the previous-only tail term. This is a deliberate implementation detail of the current code.\par
\noindent - Averages over particles include zero values for particles without collision or without valid denominator in the current interval.\par
\noindent - The CSV expansion helper writes scalar fields such as time only once at the top of an extracted block, leaving the remaining rows blank in that scalar column.\par
\noindent - If a denominator is zero, if no valid term exists, or if a feature flag is disabled in \code{Stress_settings.txt}, the corresponding exported values remain 0.\par

\end{document}
